% Lösungen Einheit 4 von 4 – Einheit 4 (zusammenbau v0.1)
\einheitenkopf{Einheit 4 von 4 · Einheit 4}

\erg{17}{a) als Summe: Quader plus Dachprisma \quad b) Pyramide $PQUT$: $V = \tfrac{1}{3} \cdot 18 \cdot 8 = 48$; über der Deckfläche ragt die ähnliche Pyramide mit den Kanten 3, 3 und 4 heraus: $V = \tfrac{1}{3} \cdot 4{,}5 \cdot 4 = 6$; Teilkörper $= 48 - 6 = 42$ \quad c) In der Ebene gilt $z = 2y$; der kleinere Teilkörper ist ein Prisma der Länge 6 mit dreieckigem Querschnitt (Katheten 3 und 6): $V = \tfrac{1}{2} \cdot 3 \cdot 6 \cdot 6 = 54$ \quad d) Die verlängerten Seitenkanten treffen sich in der Spitze der ganzen Pyramide; die Deckfläche halbiert die Seite, also liegt die Spitze 6\,m über der Deckfläche und 12\,m über der Grundfläche. Erster Summand: ganze Pyramide (Grundfläche $36\,\mathrm{m}^2$, Höhe 12\,m), zweiter: abgeschnittene Spitze (Grundfläche $9\,\mathrm{m}^2$, Höhe 6\,m); $V = 144 - 18 = 126\,\mathrm{m}^3$ \quad e) $V_Q = G \cdot h$, $V_P = \tfrac{1}{3} \cdot G \cdot h$, $V_Q : V_P = 3$; Zuwachs $= 200\,\%$ \quad f) Kegel mit $r = 6$, $h = 8$, Mantellinie $s = 10$; $O = \pi \cdot 36 + \pi \cdot 6 \cdot 10 = 96\pi \approx 301{,}59$ \quad g) $10\,\mathrm{l} = 10\,000\,\mathrm{cm}^3$, $h = 10\,000 : (\pi \cdot 10^2) \approx 31{,}83\,\mathrm{cm}$ – höher als 30\,cm, sie passt nicht \quad h) Das Teilstück mit der Ecke P ist eine zu $OPQR$ ähnliche Pyramide mit dem Streckfaktor $\tfrac{6 - t}{6}$; Volumenhalbierung: $\left(\tfrac{6 - t}{6}\right)^3 = \tfrac{1}{2}$, $t = 6 - \tfrac{6}{\sqrt[3]{2}} \approx 1{,}24$}
\erg{18}{a) Paul zieht den Teil der Pyramide nicht ab, der über den Quader hinausragt. Diese Spitze ist ähnlich mit den Kanten 4, 4 und 6: $V = \tfrac{1}{3} \cdot 8 \cdot 6 = 16$; Teilkörper $= 54 - 16 = 38$ \quad b) Die Teilpyramide ist ähnlich: Grundfläche wächst mit dem Quadrat, die Höhe mit dem Faktor selbst, das Volumen als Produkt daher mit dem Kubus. \quad c) Gerade von $(0 | 40)$ bis $(6 | 10)$ \quad d) $V = \tfrac{1}{3} \cdot 1{,}44 \cdot 18 - \tfrac{1}{3} \cdot 0{,}64 \cdot 12 = 8{,}64 - 2{,}56 = 6{,}08\,\mathrm{m}^3$, Masse $= 6{,}08 \cdot 2{,}4 \approx 14{,}59\,\mathrm{t}$}
